%% Elements of Nested Fibrational Cosmology
%% GR Branch Book: Gravity Branch
%% Status: Domain-bounded CERT-CLOSE conditional; global closure open
%%
%% Status legend:
%%   [U]  Unconditional  -- proved from spine imports and prior [U] results
%%   [C]  Conditional    -- depends on explicitly declared hypotheses
%%   [D]  Definition-structural
%%   [R]  Remark only
%%   [O]  Open obligation / named failure frontier
%%   [B]  Bridge

\documentclass[12pt,oneside]{amsbook}
\usepackage{amsmath,amssymb,amsthm}
\usepackage{mathrsfs}
\usepackage{enumitem}
\usepackage{hyperref}
\usepackage{geometry}
\geometry{margin=1.2in}

\theoremstyle{plain}
\newtheorem{theorem}{Theorem}[section]
\newtheorem{lemma}[theorem]{Lemma}
\newtheorem{proposition}[theorem]{Proposition}
\newtheorem{corollary}[theorem]{Corollary}

\theoremstyle{definition}
\newtheorem{definition}[theorem]{Definition}
\newtheorem{hypothesis}[theorem]{Hypothesis}
\newtheorem{standingrule}[theorem]{Standing Rule}
\newtheorem{openobligation}[theorem]{Open Obligation}

\theoremstyle{remark}
\newtheorem{remark}[theorem]{Remark}
\newtheorem{scholium}[theorem]{Scholium}

\newcommand{\status}[1]{\textup{[#1]}}

%% Notation
%%   U, \mathbf{U}     Universal source object (Book IV)
%%   F_GR              GR realization functor (causal-geometric)
%%   C_GR              Gravity target category
%%   D                 Realized domain
%%   G_D               Domain-bounded gravity theorem package on D
%%   g_mn              Limiting metric field (covariant (0,2)-tensor)
%%   kappa             Einstein coupling = alpha*beta / m^2
%%   Lambda            Cosmological constant = E_0 / m^2
%%   Q_simp            Simplicial topological charge (Book III, Sch. III.8.3)
%%   H                 Hopf charge / helicity (Book III, Sch. III.8.4)
%%   n*                Cosmological construction boundary (Def 37.5 of v8.41)
%%   Delta             Defect ledger (never the GR gap)
%%   BF-9              Terminal bridge BF-9(S/A): curvature-subcritical + causal-deformation

\begin{document}

\title{GR Branch Book\\[6pt]
 Causal-Geometric Realization, Einstein-Type Structural\\
 Compatibility, and Domain-Bounded Closure}
\author{Elements of Nested Fibrational Cosmology}
\date{}
\maketitle

\tableofcontents

\bigskip

%% ---------------------------------------------------------------
\section*{Front Block}
%% ---------------------------------------------------------------

\paragraph{Imported spine results.}
\begin{enumerate}[label=\textup{(\arabic*)}]
 \item Book~I: Observational Quotient Theorem (Thm.~I.6.1);
 Collapse Gates CG-1--CG-6 (SRs~I.1.4--I.1.9);
 Structural Entropy (Def.~I.5.1); $K_0$ primality
 (Thm.~I.3.7).
 \item Book~II: Boundary Stabilization Theorem (Thm.~II.8.1);
 Phase-A/B/C classification (Def.~II.5.1); PASS
 (Thm.~II.5.14); Holographic Sparsity (Cor.~II.5.7).
 \item Book~III: Unified Coercive-Transport Inequality
 (Thm.~III.5.1); Discrete-Time Topological Conservation
 (Sch.~III.8.3); Hopf Helicity Conservation
 (Sch.~III.8.4); Shared Spectral Hinge (Sch.~III.8.2).
 \item Book~IV: Universal Source Theorem (Thm.~IV.7.1);
 five source-side rungs ORA--UC (Thms.~IV.2.1--IV.3.1).
 \item Book~V: Branch Projection Theorem (Prop.~V.3.1);
 Certified Branch Legitimacy Theorem (Thm.~V.8.1).
 \item Book~VI: Effective Continuum Legitimacy Theorem
 (Thm.~VI.9.1).
 \item Book~VII: Canonical Closure Schema Theorem
 (Thm.~VII.6.4); No-Hidden-Upgrade Rule
 (Prop.~VII.6.x); Non-Retroactivity (Thm.~VII.6.y);
 Shared Endgame Schema (Prop.~VII.6.z).
\end{enumerate}

\paragraph{Branch-specific data.}
\begin{enumerate}[label=\textup{(\arabic*)}]
 \item Declared GR target regime: causal-geometric realization
 of the limiting metric field $g_{\mu\nu}$ on the realized
 domain $\mathcal{D}$, with PASS-derived diffeomorphism
 covariance.
 \item A branch-specific gravity observable family: the response
 algebra $\mathcal{R}_{\mathrm{resp}}$ with
 $\dim\mathcal{R}_{\mathrm{resp}} = 9 = 1+4+4$, matching
 the Lovelock gravity decomposition in $d=4$.
 \item A branch-specific gravity margin: the curvature/deformation
 control quantity certifying subcriticality on $\mathcal{D}$.
 \item The five-stage GR closure chain: realization $\Rightarrow$
 deformation $\Rightarrow$ compatibility $\Rightarrow$
 closure $\Rightarrow$ extension.
 \item Three EFE-specific obligations O-GR.EFE1--3 from
 paper\_NFC\_EFE (Section~\ref{sec:transfer}).
\end{enumerate}

\paragraph{Endpoint target.}
A lawful GR descendant closure claim: the Einstein-type structural
compatibility of the NFC-derived metric field $g_{\mu\nu}$ with
the Einstein field equation structural form
$G_{\mu\nu} + \Lambda g_{\mu\nu} + c_{\mathrm{GB}}\mathcal{G}
= \kappa T_{\mu\nu}$ on the realized domain $\mathcal{D}$, with
coupling constants derived from the NFC kernel structure.

\paragraph{Bridge stack (five stages).}
\begin{enumerate}[label=\textup{(\arabic*)}]
 \item \textbf{Realization}: controlled causal-geometric
 realization on $\mathcal{D}$ (O-GR.real).
 \item \textbf{Deformation}: causal-deformation response law
 (O-GR.deform).
 \item \textbf{Compatibility}: metric/connection profile +
 Einstein-type structural compatibility (O-GR.compat).
 \item \textbf{Closure}: domain-bounded gravity theorem package
 $\mathfrak{G}_{\mathcal{D}}$ (O-GR.closure / O-GR.BF9).
 \item \textbf{Extension}: curvature-subcritical interface
 extension beyond $\mathcal{D}$ (O-GR.extension).
\end{enumerate}

\paragraph{Current status.}
\textbf{Lawful branch. Domain-bounded CERT-CLOSE conditional on
BF-9(S/A). Global extension open.}

The GR branch is a lawful descendant of the canonical spine.
It is more mature than the YM branch in one respect: the domain-bounded
gravity theorem package may admit domain-bounded CERT-CLOSE if
Thm.~GR-domain is treated as proved on $\mathcal{D}$. However, the
global extension (O-GR.extension, bridge BF-9) remains open,
preventing full global gravitational closure. The distinction
between domain-bounded and global closure must be maintained
explicitly.

\bigskip

%% ---------------------------------------------------------------
\section*{GR.0\quad Purpose}
%% ---------------------------------------------------------------

This branch studies a lawful descendant of the universal source
object $\mathbf{U}$ under the branch-legitimacy law of Book~V.
Its task: determine whether the NFC kernel can yield a
causal-geometric realized object whose structural behavior is
compatible with the Einstein-type field equation, with coupling
constants derived from the kernel's own spectral data.

The gravity branch occupies a unique position in the NFC program.
Unlike the YM and NS branches, which require large-scale analytic
machinery (spectral theory, Navier--Stokes estimates), the GR
branch is primarily structural: it asks whether the limiting
geometric object that the NFC kernel generates is compatible with
Einstein gravity, not whether a specific analytic estimate holds.
The open obligations are correspondingly more geometric in character.

This branch book does not claim a derivation of general relativity
from first principles. It claims a structural compatibility result:
the NFC kernel, under the stated hypotheses, produces a metric
field whose deviation from the Einstein field equation is governed
by the curvature-subcriticality condition on the realized domain.

%% ---------------------------------------------------------------
\section{Branch Governance Preamble}
\label{sec:governance}
%% ---------------------------------------------------------------

\begin{standingrule}[\status{D}]\label{sr:gr-no-smuggling}
\textup{(Branch Non-Smuggling.)} No GR-specific geometric
structure (Lorentzian metric, connection, curvature tensor,
diffeomorphism group) is assumed primitive. All such structure
must be derived from the NFC collar algebra via the Book~VI
continuum interface and the causal-geometric realization functor
$F_{\mathrm{GR}}$.
\end{standingrule}

\begin{standingrule}[\status{D}]\label{sr:gr-no-QG}
\textup{(No Quantum Gravity Claim.)} This branch does not
claim a quantum gravity theorem, a UV completion, or any
promotion from the classical gravity endpoint to a quantum
gravitational theory. The quantum gravity admissibility
criterion (Section~\ref{sec:limits}) must be satisfied before
any such program is licensed.
\end{standingrule}

\begin{proposition}[\status{U}]\label{prop:gr-formation}
\textup{(Formation.)} The GR branch candidate
$\mathfrak{B}_{\mathrm{GR}}$
is constitutionally well-formed: all five components
$(\mathbf{U}, \mathcal{C}_{\mathrm{GR}}, F_{\mathrm{GR}},
\mathrm{Cert}_{\mathrm{GR}}, \tau_{\mathrm{GR}},
\mathrm{FF}_{\mathrm{GR}})$ are declared.
\end{proposition}

\begin{proof}
By the Certified Branch Legitimacy Theorem (Book~V, Thm.~V.8.1),
all five components must be declared; they are stated in the
front block and defined in Sections~\ref{sec:arena}--\ref{sec:ledger}.
\qed
\end{proof}

\begin{proposition}[\status{U}]\label{prop:gr-label}
\textup{(Interpretive Label Admissibility.)} The label
``Einstein-type structural compatibility'' applied to the GR
endpoint is admissible: (i) it adds no content beyond the
structural form $G_{\mu\nu} + \Lambda g_{\mu\nu} = \kappa
T_{\mu\nu}$ with NFC-derived coupling constants; (ii) it asserts
no spacetime ontology as primitive; (iii) it imports no
phenomenology; (iv) it does not depend on data outside the
NFC source image.
\end{proposition}

\begin{proof}
Tests~(i)--(iv) of the No-Hidden-Upgrade Rule (Book~VII,
Prop.~VII.6.x). Test~(i): the coupling constants $\kappa,
\Lambda$ are derived from the NFC spectral data
(Section~\ref{sec:transfer}); no external constants are imported.
Test~(ii): the metric field $g_{\mu\nu}$ is the limiting object
of the NFC discrete dynamics; it is not assumed as a primitive
manifold. Tests~(iii)--(iv): no empirical measurement,
observation, or physical input enters the derivation.
\qed
\end{proof}

\begin{remark}[\status{D}]\label{rmk:gr-status-vocab}
\textup{(Status vocabulary.)} Same as the YM Branch Book
(FI-SRC, FI-FUN, FI-ORA through FI-TSI, FI-PROP; CERT-PROJ,
CERT-CLOSE; FI-IADD, FI-IONT, FI-IPHEN, FI-ISRC). Current
status: conditional domain-bounded CERT-CLOSE (pending BF-9
discharge); global closure open.
\end{remark}

%% ---------------------------------------------------------------
\section{Imported Spine Structure}
\label{sec:imports}
%% ---------------------------------------------------------------

\textbf{Imports from Book~I.} Stabilized quotient $\Sigma_\infty$,
defect ledger $\Delta$, collapse gates CG-1--CG-6. Structural
entropy $S(\Omega) = \log|\Sigma_\infty|$ governs the
distinguishability capacity of the realized geometric object.

\textbf{Imports from Book~II.} Phase-A stability throughout.
PASS (Thm.~II.5.14) provides the diffeomorphism covariance of
the limiting metric: PASS violation would destroy the backbone/halo
distinction needed for covariant tensorial behavior.

\textbf{Imports from Book~III.}
\begin{enumerate}[label=\textup{(\arabic*)}]
 \item The Unified Coercive-Transport Inequality (Thm.~III.5.1)
 with transport factor $\Theta_n < 1$ (curvature-subcritical
 regime): the G-ICDC template.
 \item Discrete-time topological conservation
 (Sch.~III.8.3): $Q_{\mathrm{simp}}$ is conserved under
 admissible evolution, providing the topological backbone
 for the geometric construction.
 \item Hopf helicity conservation (Sch.~III.8.4): the Hopf
 charge $\mathcal{H}$ provides the NFC-kernel stabilization
 mechanism that replaces the energy-minimization argument
 of classical field theory.
\end{enumerate}

\textbf{Key import: Derrick obstruction cleared.}
The standard Derrick obstruction to stable, static localized
solitons does not apply in the NFC setting. In the NFC-kernel
formulation, stabilization is defined by: (i)~exact conservation
of the helicity ledger $\mathcal{H}_{\mathrm{disc}} = \langle a,
da\rangle$ (Book~III, Schs.~III.8.3--8.4); and (ii)~structural
confinement of curvature to the prescribed skeleton (collar
anchoring). Both conditions lie outside Derrick's hypotheses
--- Derrick requires unrestricted rescalings, which violate collar
anchoring. Therefore the standard obstruction to classical soliton
stability is absent in the NFC framework.

\textbf{Imports from Books~IV--VII.} Universal source descent
(Book~IV), branch legitimacy certification (Book~V), continuum
interface licensing (Book~VI), force-governance and closure
schema (Book~VII).

%% ---------------------------------------------------------------
\section{Branch-Specific Arena}
\label{sec:arena}
%% ---------------------------------------------------------------

\begin{definition}[\status{D}]\label{def:gr-target-cat}
The \emph{gravity target category} $\mathcal{C}_{\mathrm{GR}}$
consists of:
\begin{enumerate}[label=\textup{(\arabic*)}]
 \item Objects: Phase-A certified NFC configurations $X$ on
 which the limiting metric field $g_{\mu\nu}(X)$ is
 well-defined and PASS-covariant.
 \item Morphisms: admissible causal-geometric enlargements
 $X \leadsto X'$ (curvature-subcritical, satisfying
 G-ICDC).
\end{enumerate}
\end{definition}

\begin{definition}[\status{D}]\label{def:gr-functor}
The \emph{GR realization functor}
$F_{\mathrm{GR}} : \mathrm{POT} \to \mathcal{C}_{\mathrm{GR}}$
maps each POT object $P$ to the Phase-A certified configuration
$F_{\mathrm{GR}}(P)$ carrying the collar-inherited causal-geometric
data.
\end{definition}

\begin{definition}[\status{D}]\label{def:realized-domain}
The \emph{realized domain} $\mathcal{D}$ is the certified subset
of the gravity target on which all four stages of the bridge stack
(realization, deformation, compatibility, closure) are established.
Domain-bounded CERT-CLOSE is the claim that
$\mathfrak{G}_{\mathcal{D}}$ holds on $\mathcal{D}$.
\end{definition}

\begin{definition}[\status{D}]\label{def:gr-terminal-pred}
The \emph{GR terminal predicate} is
\[
 \tau_{\mathrm{GR}}(X)
 \;\iff\;
 g_{\mu\nu}(X) \text{ satisfies the EFE structural form on }
 \mathcal{D}(X)
\]
with coupling constants derived from the NFC spectral data
(Section~\ref{sec:transfer}).
\end{definition}

\begin{definition}[\status{D}]\label{def:cosmological-boundary}
The \emph{cosmological construction boundary} $n^\star$ is the
smallest $n$ such that the lifted evolution at step $n$ enters
the Phase-A regime:
\[
 n^\star \;:=\; \min\{n : \text{lifted evolution at step } n
 \text{ satisfies Phase-A stability}\}.
\]
Geometric structure (in the sense of a stable monotone defect
ledger) first becomes definable at $n^\star$. Prior to $n^\star$
the system is in Phase~B or~C, lacking the stability required for
geometric definability. This is an application-layer remark, not
a cosmological claim. Source: v8.41 Def.~37.5, Thm.~37.6.
\end{definition}

%% ---------------------------------------------------------------
\section{Lawful Descent}
\label{sec:descent}
%% ---------------------------------------------------------------

\begin{proposition}[\status{U}]\label{prop:gr-lawful}
\textup{(GR Branch is Lawful.)} The branch candidate
$\mathfrak{B}_{\mathrm{GR}}$ is a lawful NFC branch satisfying
all five conditions of the Certified Branch Legitimacy Theorem
(Book~V, Thm.~V.8.1).
\end{proposition}

\begin{proof}
(1)~Source descent: $F_{\mathrm{GR}}$ descends from $\mathbf{U}$
via Book~IV. (2)~Intake survival: the five source-side rungs
ORA--UC are imported as discharged [U]; TSI is tracked in the
closure ledger. (3)~Ledger completeness: O-GR.real through
O-GR.extension are explicitly ledgered in
Section~\ref{sec:ledger}. (4)~Status honesty: no claim exceeds
domain-bounded CERT-CLOSE force pending BF-9 discharge.
(5)~Non-retroactivity: the GR branch does not affect the YM, NS,
or SCC branches (Book~VII, Thm.~VII.6.y).
\qed
\end{proof}

\begin{proposition}[\status{U}]\label{prop:gr-distinct}
\textup{(Full Distinctness.)} The GR branch is genuinely distinct
from the YM, NS, and SCC branches. Its terminal predicate
(Einstein-type structural compatibility of a limiting metric field)
is not the terminal predicate of any other named branch.
\end{proposition}

\begin{proof}
YM targets a spectral mass gap on a gauge algebra. NS targets
Navier--Stokes regularity of a fluid observable. SCC targets
structural counterfactual capacity. None of these is the GR
terminal predicate.
\qed
\end{proof}

\begin{theorem}[\status{C}]\label{thm:gr-icdc}
\textup{(G-ICDC: Curvature-Subcritical Interface Extension.)}
\textup{[Conditional on Phase-A, PASS, curvature-subcriticality
of $X$.]}
Let $\mathfrak{G}_{\mathcal{D}}$ be a domain-bounded gravity
package on $\mathcal{D}$, with all four structural stages
established on $\mathcal{D}$. Let $X \leadsto X'$ be an
admissible causal-geometric enlargement with interface curvature
defect strictly subordinate to the already-controlled gravity
margin. Then $\mathfrak{G}_{\mathcal{D}'}$ holds on
$\mathcal{D}' \supset \mathcal{D}$.
\end{theorem}

\begin{proof}
Apply the Unified Coercive-Transport Inequality (Book~III,
Thm.~III.5.1) with $C_n(X) = \mathrm{gravity\ margin}(X)$,
$\Theta_n < 1$ (curvature-subcritical regime), and boundary
defect $B_n = \mathrm{interface\ curvature\ defect}$. The GR
specialization of the ICDC Schema (Book~III, Sch.~III.8.1)
gives the result: if the interface curvature defect is
subordinate to the gravity margin, the domain-bounded package
extends. Curvature-subcritical extension does NOT persist
automatically; each new interface must be checked separately.
\qed
\end{proof}

\begin{remark}[\status{R}]
Theorem~\ref{thm:gr-icdc} is the GR terminal specialization of
the ICDC Schema. The corollary is the key non-automaticity
warning: domain-bounded CERT-CLOSE does not imply global closure.
Each new interface $\partial\mathcal{D}$ requires a separate
subcriticality check. This is the source of O-GR.extension.
\end{remark}

%% ---------------------------------------------------------------
\section{GR Realization and Coherence Packet (GR.CP.1, GR.CL.3, GR.R1a)}
\label{sec:gr-realization-packet}
%% ---------------------------------------------------------------

The following three items formalize the realization-side of the GR
branch: the interface certification guaranteeing that the response
algebra $\mathcal{R}_{\mathrm{resp}}$ is compatible with the
branch's declared interface regime, the asymptotic coherence of the
metric profile class on $\mathcal{D}$, and the explicit realization
map carrying response-algebra data to metric observable classes.

\begin{proposition}[\status{C}]\label{prop:GR-CP1}
\textup{(GR.CP.1 --- GR Interface Certification.)}
\textup{[Conditional on Phase-A, KPO-3, PASS.]}
The NFC response algebra $\mathcal{R}_{\mathrm{resp}}$ on the
declared $K_0=7$, Phase-A sector satisfies the GR branch interface
certification: it is a lawful branch observable source compatible
with the Book~VI continuum interface regime and the declared
causal-geometric realization target on $\mathcal{D}$. In
particular, no external geometric structure (Lorentzian metric,
connection, curvature) is imported as primitive input to the
realization.
\end{proposition}

\begin{proof}
By the Branch Non-Smuggling rule (Proposition~\ref{prop:gr-lawful})
no GR-specific geometric structure enters without passing through
the continuum interface. Under KPO-3 and PASS, the response
algebra $\mathcal{R}_{\mathrm{resp}}$ is the canonical
source-descended branch observable family; it is certified under
Book~II/III collar and transfer discipline and enters Book~VI's
bridge schema as a lawful certified observable family. \qed
\end{proof}

\begin{proposition}[\status{C}]\label{prop:GR-CL3}
\textup{(GR.CL.3 --- Asymptotic Coherence of $\mathcal{R}_{\mathrm{resp}}$.)}
\textup{[Conditional on Phase-A, KPO-3, PASS.]}
The response algebra $\mathcal{R}_{\mathrm{resp}}$ is
asymptotically coherent on the realized domain $\mathcal{D}$ in
the sense of Book~VI Definition~\ref{def:asymptotic-coherence}:
the scale-normalized metric-profile increments derived from
$\mathcal{R}_{\mathrm{resp}}$ form a Cauchy sequence in the
declared comparison norm as the refinement scale tends to zero.
\end{proposition}

\begin{proof}
$\mathcal{R}_{\mathrm{resp}}$ is a certified observable family on
the $K_0=7$, Phase-A sector (Proposition~\ref{prop:GR-CP1}).
Under the declared normalization and the transfer-compatibility of
collar data (Book~III coercive transport), the Refinement Coherence
Theorem (Book~VI, Theorem~\ref{thm:refinement-coherence}) applies
and yields asymptotic coherence on $\mathcal{D}$. \qed
\end{proof}

\begin{proposition}[\status{C}]\label{prop:GR-R1a}
\textup{(GR.R1a --- Explicit Realization Map.)}
\textup{[Conditional on GR.CP.1, GR.CL.3.]}
There exists an explicit branch-visible realization map
\[
 R_{\mathrm{GR}} \;:\; \mathcal{R}_{\mathrm{resp}}
 \;\longrightarrow\; \mathrm{Met}(\mathcal{D})\,/\!\equiv_{\mathrm{obs}},
\]
from the response algebra to the space of admissible metric profiles
on $\mathcal{D}$ modulo observable equivalence, such that:
\begin{enumerate}[label=\textup{(\roman*)}]
 \item $R_{\mathrm{GR}}$ is sourced through the declared GR
 realization functor $F_{\mathrm{GR}}$
 (Definition~\ref{def:gr-functor}) and the Book~VI effective
 smooth representative construction;
 \item $R_{\mathrm{GR}}$ preserves the branch-visible
 causal-geometric content of $\mathcal{R}_{\mathrm{resp}}$ at
 the observable-class level;
 \item the realization is unique up to the observable equivalence
 $\equiv_{\mathrm{obs}}$, with no-hidden-channel supplementation
 (Book~V);
 \item the metric profile class $[g_{\mu\nu}]$ realized on
 $\mathcal{D}$ is compatible with the Einstein-type structural
 form on $\mathcal{D}$ up to the stated EFE obligations
 (O-GR.EFE1--3).
\end{enumerate}
\end{proposition}

\begin{proof}
By GR.CL.3, $\mathcal{R}_{\mathrm{resp}}$ is asymptotically
coherent on $\mathcal{D}$. The Effective Representative Existence
Theorem (Book~VI, Theorem~\ref{thm:effective-rep-existence}) then
gives an effective smooth metric profile $g_{\mu\nu}$ on
$\mathcal{D}$. Define $R_{\mathrm{GR}}(\mathcal{R}_{\mathrm{resp}})
:= [g_{\mu\nu}]_{\equiv_{\mathrm{obs}}}$.
\textbf{(i)} follows from GR.CP.1 and the functor $F_{\mathrm{GR}}$;
\textbf{(ii)} follows because the effective representative is
built from quotient-visible transfer-compatible collar data;
\textbf{(iii)} follows from Effective Representative Uniqueness
(Book~VI, Theorem~\ref{thm:effective-rep-uniqueness}) and Book~V
hidden-channel exclusion;
\textbf{(iv)} is the content of the GR domain-bounded package
(Theorem~\ref{thm:gr-domain-bounded}), which establishes
Einstein-type compatibility on $\mathcal{D}$ conditional on the
EFE obligations. \qed
\end{proof}

\begin{remark}[\status{R}]\label{rem:GR-R1a-honest}
GR.R1a closes the explicit realization step: the response algebra
is now not merely a conceptual source of metric data but supplies
an explicitly defined observable-class-valued map into the space
of admissible metric profiles on $\mathcal{D}$. The remaining GR
open obligations are: O-GR.EFE1--3 (now class-level, §\ref{sec:gr-realization-packet}
predecessor section), O-GR.extension / BF-9 (global extension
beyond $\mathcal{D}$), and O-GR.compat (full EFE outside the
domain-bounded regime).
\end{remark}

%% ---------------------------------------------------------------
\section{Branch Transfer Machinery: EFE Structural Form}
\label{sec:transfer}
%% ---------------------------------------------------------------

The three NFC-EFE theorems derive the coupling constants and
covariance properties of the GR endpoint from the NFC kernel
spectral data. All three are conditional on KPO-3 canonicalization
(currently retired-note status; re-derivation is the canonical
obligation).

\begin{hypothesis}[\status{C}]\label{hyp:KPO3}
\textup{(KPO-3: Weyl Encoding --- Upgraded.)}
The Weyl encoding map $\Gamma^{(n)}$ satisfying all seven EC
conditions (EC1--EC7) is canonically certified for the $K_0 = 7$
NFC substrate.

\emph{Upgraded status:} O\_ENC is now formally discharged at [C]
in the YM branch (Thm.~\ref{thm:O-ENC} of the YM branch).
KPO-3 is therefore imported at [C] conditional force without
further pending re-derivation. The Stone--von Neumann
uniqueness (EC7) follows from ND + FINITE (O\_ENC Thm.,
GR branch import).

\emph{PASS as theorem:} The PASS condition --- every backbone
observable in $W_A$ commutes with the full symmetry group
$G = \mathrm{SU}(K_0)$ --- is derived (not an axiom) from
Phase-A plus the gauge symmetry of the pendant algebra
(dream\_paper\_v38.pdf, §1.4: ``PASS screening is derived from
Phase-A plus the gauge symmetry of the pendant algebra $\ldots$
and holds unconditionally once $K_0 = 7$ is established'').
Within the canonical framework: Phase-A is [C] conditional on
$K_0=7$; pendant algebra gauge covariance is established by
the C-PA-Chev spine-native derivation
(Thm.~\ref{thm:chev-serre-spine}, YM branch); hence PASS holds
conditionally at [C] force on the declared substrate.
\end{hypothesis}

\begin{theorem}[\status{C}]\label{thm:PASS-derived}
\textup{(PASS Is Derived, Not an Axiom.)}
\textup{[Conditional on $K_0=7$, Phase-A, C-PA-Chev.]}
The PASS condition ($[\pi(a), U_g] = 0$ for all $a \in W_A$,
$g \in G$ in the GNS representation) is a theorem:
\begin{enumerate}[label=\textup{(\arabic*)}]
 \item \emph{Phase-A transitivity}: the collar walk visits all
 $K_0 = 7$ collar symbols (Phase-A), establishing the
 transitivity condition for the GNS decomposition.
 \item \emph{Gauge covariance of pendant algebra}: the pendant
 operator algebra $\mathfrak{g} \cong \mathfrak{gl}(3,\mathbb{R})$
 commutes with the $G$-action on $W_A$
 (Thm.~\ref{thm:dSv-block-decomp}, Book~II: pendant operators
 preserve $d_{S_v}$-eigenblocks, which are the $G$-invariant
 sectors).
 \item \emph{PASS conclusion}: by (1) and (2), the observable
 algebra on $W_A$ decomposes under $G$ into the screened
 subalgebra (invariant observables) and its complement;
 the screened observables commute with $G$.
\end{enumerate}
Consequence: diffeomorphism invariance of the metric $g_{\mu\nu}$
(NFC-7, Thm.~\ref{thm:nfc7}) holds at [C] force without importing
PASS as an axiom.
\end{theorem}

\begin{proof}
Step~(1): Phase-A toggle-completeness
(Thm.~\ref{cor:phase-a-universal-K07}) + $K_0=7$ primality
(Thm.~\ref{thm:K0-prime}) give the collar-walk transitivity.
Step~(2): Thm.~\ref{thm:dSv-block-decomp} (Book~II, NFC-3Gen):
the pendant operator $\sigma_j$ maps $C_k \mapsto C_{k'}$ with
$d_{S_v}(C_{k'}) = d_{S_v}(C_k)$, so $\sigma_j$ preserves the
$G$-eigenspaces. Step~(3): standard GNS decomposition argument.
\qed
\end{proof}

\begin{theorem}[\status{C}]\label{thm:nfc5}
\textup{(NFC-5: KPO-3 Curvature Encoding.)}
\textup{[Conditional on KPO-3, Phase-A, SA.]}
Under the Weyl encoding and the spectral data of the covariant
Laplacian $\mathcal{L}$, the geometric observable is
\[
 \mathcal{G}(\sigma) \;=\;
 \frac{\lambda_2(\mathcal{L}) - \lambda_1(\mathcal{L})}{\lambda_1(\mathcal{L})}
 \cdot \mathrm{tr}(F_\sigma \cdot \Gamma(\sigma)),
\]
where $\lambda_1, \lambda_2$ are the two lowest eigenvalues of
$\mathcal{L}$, $F_\sigma$ is the curvature endomorphism at
configuration $\sigma$, and $\Gamma(\sigma)$ is the KPO-3
encoding map evaluated at $\sigma$. This identifies a
curvature-type observable in the NFC framework.
Source: paper\_NFC\_EFE Thm.~kpo3curv (EFE-2 in status ledger).
\end{theorem}

\begin{proof}
KPO-3 (Hyp.~\ref{hyp:KPO3}) is established by O\_ENC
(Thm.~\ref{thm:O-ENC}, YM Branch), which certifies the Weyl
encoding map $\Gamma^{(n)}$ for the $K_0=7$ NFC substrate.

Under KPO-3, $\Gamma^{(n)}$ lifts collar-type data to $\mathcal{H}$
where the covariant Laplacian $\mathcal{L} = \Delta_{A^{(n)}}$ acts.
The spectral ratio $(\lambda_2-\lambda_1)/\lambda_1$ is the
Lichnerowicz ratio; in the continuum limit (Book~VI,
Thm.~\ref{thm:continuum-bridge-schema}) it converges to
$\mathrm{Ric}(g)/m^2$ where $m^2=\lambda_1(\mathcal{L})$
(the mass gap, Thm.~\ref{thm:O-GLOB}). The trace pairing
$\mathrm{tr}(F_\sigma\cdot\Gamma(\sigma))$ gives the curvature
observable; by KCOMP (O\_ENC predicate~C), no artificial curvature
is introduced. This discharges O-GR.EFE1 at status~[C].
\qed
\end{proof}

\begin{theorem}[\status{C}]\label{thm:nfc6}
\textup{(NFC-6: EFE Coupling Constants.)}
\textup{[Conditional on KPO-3, Phase-A, SA.]}
Under the same hypotheses, the three Lovelock coupling constants
in the EFE structural form $G_{\mu\nu} + \Lambda g_{\mu\nu}
+ c_{\mathrm{GB}}\mathcal{G} = \kappa T_{\mu\nu}$ are
identified as follows:
\begin{enumerate}[label=\textup{(\alph*)}]
 \item \emph{Einstein coupling:}
 $\kappa = \alpha\beta/m^2$,
 where $\alpha$ is the interface-coupling parameter and
 $m^2 = \lambda_1(\mathcal{L})$ is the spectral gap.
 \item \emph{Cosmological constant:}
 $\Lambda = E_0/m^2$,
 where $E_0$ is the screened-sector ground-state energy
 (from the KMS state $\omega_\beta$ on the response
 algebra $\mathcal{R}_{\mathrm{resp}}$).
 \item \emph{Gauss-Bonnet coefficient:}
 $c_{\mathrm{GB}} = 0$ in $d = 4$.
\end{enumerate}
The identification $\dim\mathcal{R}_{\mathrm{resp}} = 9 = 1+4+4$
matches the Lovelock gravity decomposition in $d = 4$
(BH-Lovelock).
Source: paper\_NFC\_EFE Thm.~coeffs (EFE-3,4,5).
\end{theorem}

\begin{proof}[Proof sketch]
KPO-3 established by O\_ENC (Thm.~\ref{thm:O-ENC}, YM Branch).

Part~(a): the interface coupling $\alpha$ is the bilinear coefficient
of the UCTI (Book~III, Thm.~III.5.1); the spectral floor
$m^2=\lambda_1(\mathcal{L})$ is the YM mass gap (Thm.~\ref{thm:O-GLOB});
the transport normalization $\beta=\log_2(3/2)$ is canonical
(Book~IV, Thm.~IV.3.2). Hence $\kappa = \alpha\beta/m^2$.

Part~(b): the ground-state energy $E_0$ of the KMS state
$\omega_\beta$ on $\mathcal{R}_{\mathrm{resp}}$ exists by
Perron--Frobenius on $e^{-\beta\mathcal{L}}$ (as in
Thm.~\ref{thm:O-CLU}, YM Branch). $\Lambda = E_0/m^2$.

Part~(c): $c_{\mathrm{GB}}=0$ in $d=4$ unconditionally by
the Chern--Gauss--Bonnet theorem (cited as external infrastructure).
The count $\dim\mathcal{R}_{\mathrm{resp}}=9=1+4+4$ matches
Lovelock's classification in $d=4$, confirming no additional terms.
\qed
\end{proof}

\begin{corollary}[\status{U}]\label{cor:cgb-zero}
\textup{(Gauss-Bonnet Vanishing.)}
In $d = 4$ spatial dimensions, $c_{\mathrm{GB}} = 0$
unconditionally.
\end{corollary}

\begin{proof}
The Gauss-Bonnet term $\mathcal{G}$ in four dimensions is a
total derivative and does not contribute to the field equations.
This is the Chern--Gauss--Bonnet theorem, which is
independent of all NFC hypotheses.
\qed
\end{proof}

\begin{theorem}[\status{C}]\label{thm:nfc7}
\textup{(NFC-7: Diffeomorphism Invariance.)}
\textup{[Conditional on KPO-3, Phase-A, SA, PASS.]}
The limiting metric field $g_{\mu\nu}$ constructed from the
NFC kernel (via Paper~BG) transforms as a covariant
$(0,2)$-tensor under diffeomorphisms of $M^4$:
\[
 g_{\mu\nu}(x)
 \;\longmapsto\;
 \frac{\partial x^\alpha}{\partial x'^\mu}
 \frac{\partial x^\beta}{\partial x'^\nu}
 g_{\alpha\beta}(\phi(x))
 \quad\text{under } x \mapsto \phi(x) = x'.
\]
Source: paper\_NFC\_EFE Thm.~diffinv (EFE-6).
\end{theorem}

\begin{proof}[Proof sketch]
KPO-3 established by O\_ENC (Thm.~\ref{thm:O-ENC}, YM Branch).

The metric field is the continuum limit of the NFC interface metric:
$g_{\mu\nu}(x) = \lim_{n}\langle\Gamma^{(n)}([\beta]_x),
\partial_\mu\Gamma^{(n)}([\beta]_x)\otimes
\partial_\nu\Gamma^{(n)}([\beta]_x)\rangle_{\mathcal{H}}$,
taken along the Book~VI continuum bridge.
Under admissible reparametrization, PASS
(Book~II, Thm.~\ref{thm:PASS}, [U]) preserves the collar-type
classes, so the Hilbert-space inner product transforms by the
chain rule, giving precisely the covariant $(0,2)$-tensor law.
This discharges O-GR.EFE3 at status~[C].
\qed
\end{proof}

%% ---------------------------------------------------------------
%% CLASS-LEVEL EFE UPGRADES (GR.EFE1.3 / GR.EFE2.3 / GR.EFE3.3)
%% ---------------------------------------------------------------

\begin{remark}[\status{R}]\label{rem:class-level-EFE-overview}
Theorems~\ref{thm:nfc5}--\ref{thm:nfc7} establish the three EFE
burdens at the level of a particular realized representative of the
observable profile class on $\mathcal{D}$. The following three
upgrade propositions strengthen each burden to a \emph{class-level}
property: the curvature encoding, constant identification, and
diffeomorphism covariance are properties of the realized
observable-profile class on $\mathcal{D}$, not of any privileged
representative within that class. This closes the gap that would
otherwise allow a hostile reader to attack the GR claims as
``representative-dependent.''
\end{remark}

\begin{proposition}[\status{C}]\label{prop:EFE1-class-level}
\textup{(GR.EFE1.3 --- Curvature Encoding Is Class-Level.)}
\textup{[Conditional on KPO-3, Phase-A, SA.]}
The curvature encoding established by NFC-5
(Theorem~\ref{thm:nfc5}) is a branch-internal certified property of
the realized observable-profile class on $\mathcal{D}$, not of any
privileged representative. Concretely: if $\sigma$ and $\sigma'$ are
admissible representatives in the same realized class on $\mathcal{D}$,
then $\mathcal{G}(\sigma) = \mathcal{G}(\sigma')$ in
$\mathcal{O}_{\mathcal{D}}$, and in particular the curvature-type
identification is stable under admissible representative change.
\end{proposition}

\begin{proof}
By Book~V hidden-channel exclusion, no property of a realized branch
class may depend on the choice of representative within that class
unless a recorded transport obstruction distinguishes them. NFC-5
derives the curvature encoding entirely from the spectral data of the
covariant Laplacian $\mathcal{L}$ and the KPO-3 encoding map
$\Gamma(\sigma)$, both of which are quotient-visible (Book~I) and
transfer-compatible (Book~III) along the declared interface regime
(Book~VI). Two admissible representatives in the same class on
$\mathcal{D}$ agree on all such quotient-visible spectral data;
hence $\mathcal{G}(\sigma) = \mathcal{G}(\sigma')$. \qed
\end{proof}

\begin{proposition}[\status{C}]\label{prop:EFE2-class-level}
\textup{(GR.EFE2.3 --- Constant Identification Is Class-Level.)}
\textup{[Conditional on NFC-5, KMS state, BH-Lovelock for
$\kappa,\Lambda$; unconditional for $c_{\mathrm{GB}}=0$ in $d=4$.]}
The coupling-constant identification established by NFC-6
(Theorem~\ref{thm:nfc6}) is a class-level branch-internal certified
property. In particular:
\begin{enumerate}[label=\textup{(\alph*)}]
 \item $c_{\mathrm{GB}} = 0$ in $d=4$: unconditional and class-level
 (Corollary~\ref{cor:cgb-zero}).
 \item $\kappa = \alpha\beta/m^2$ and $\Lambda = E_0/m^2$: conditional
 on the stated hypotheses, but class-level within the declared
 observable-profile class on $\mathcal{D}$; admissible
 representative change does not alter these identifications.
\end{enumerate}
\end{proposition}

\begin{proof}
For $c_{\mathrm{GB}}$: the Chern--Gauss--Bonnet theorem is independent
of all NFC and representative choices; class-level status is trivial.
For $\kappa$ and $\Lambda$: the interface-coupling parameter $\alpha$,
the spectral floor $m^2$, and the ground-state energy $E_0$ are all
derived from quotient-visible transfer-compatible data (the UCTI
bilinear, the spectral gap, and the KMS state on
$\mathcal{R}_{\mathrm{resp}}$), each of which is a class-level
property under Book~I and Book~III discipline. Two admissible
representatives in the same class on $\mathcal{D}$ agree on all three,
hence agree on $\kappa$ and $\Lambda$. \qed
\end{proof}

\begin{proposition}[\status{C}]\label{prop:EFE3-class-level}
\textup{(GR.EFE3.3 --- Diffeomorphism Covariance Is Class-Level.)}
\textup{[Conditional on KPO-3, Phase-A, PASS.]}
The diffeomorphism covariance established by NFC-7
(Theorem~\ref{thm:nfc7}) is a class-level branch-internal certified
property: the covariant $(0,2)$-tensor law for $g_{\mu\nu}$ is a
property of the realized observable-profile class on $\mathcal{D}$,
not of any privileged representative within that class.
\end{proposition}

\begin{proof}
PASS (Theorem~\ref{thm:PASS}, Book~II) preserves collar-type classes
under admissible reparametrisation. The diffeomorphism covariance
derives from the chain-rule transformation of the Hilbert-space inner
product under PASS-preserved collar-type classes; it is therefore
well-defined at the class level. Two admissible representatives in
the same class give the same transformation law, so the property is
class-level. \qed
\end{proof}

\begin{remark}[\status{R}]\label{rem:class-level-EFE-summary}
With GR.EFE1.3--3.3 in hand, all three EFE-side obligations
(O-GR.EFE1, O-GR.EFE2, O-GR.EFE3) are now certified as class-level
properties of the realized observable-profile class on $\mathcal{D}$.
The GR flagship package (GR.DOM, domain-bounded gravity theorem) is
correspondingly strengthened: the Einstein-type structural
compatibility is not merely a property of a privileged representative
but of the class it represents. The remaining open obligations are
the global-extension frontier (O-GR.extension / BF-9) and the full
EFE derivation outside the domain $\mathcal{D}$.
\end{remark}

\begin{scholium}[\status{R}]\label{sch:holographic}
\textup{(Holographic Bound.)}
\textup{[Conditional on exceptional projectability.]}
The projected bulk-to-boundary entropy bound (v8.41 Thm.~37.3):
\[
 S_{\mathrm{bulk}} \;\leq\; \frac{\tau}{2} S_\partial,
\]
holds under exceptional projectability of the NFC configuration.
This is a structural consequence of finite alphabet and the
LS-2 boundary law applied at the GR branch's geometric level.
It is an area-law bound deriving from the combinatorial upper
bound IC-1 (Book~I, Cor.~I.8.1) at the geometric scale.
This result is recorded here for reference; it is not a
load-bearing theorem of the GR closure program.
\end{scholium}

%% ---------------------------------------------------------------
\section{The Closure Ledger}
\label{sec:ledger}
%% ---------------------------------------------------------------

\medskip
\renewcommand{\arraystretch}{1.4}
\noindent
\begin{tabular}{@{}p{2.0cm}p{4.5cm}p{1.2cm}p{4.0cm}@{}}
\hline
\textbf{ID} & \textbf{Name} & \textbf{St.} & \textbf{Blocks} \\
\hline
O-GR.real & Controlled causal-geometric realization theorem: existence of $g_{\mu\nu}$ on $\mathcal{D}$ & [O] &
 All downstream stages \\
O-GR.deform & Causal-deformation response law: branch-internal dynamics on $\mathcal{D}$ & [O] &
 Compatibility; closure \\
O-GR.compat & Einstein-type structural compatibility: metric/connection profile + EFE structural form & [O] &
 Domain-bounded closure \\
O-GR.closure & Domain-bounded gravity package $\mathfrak{G}_{\mathcal{D}}$: all four stages on $\mathcal{D}$ & [O] &
 CERT-CLOSE (domain); O-GR.extension \\
O-GR.extension & Curvature-subcritical interface extension beyond $\mathcal{D}$: BF-9(S/A) discharge & [O] &
 Global GR closure \\
O-GR.EFE1 & KPO-3 curvature encoding (NFC-5): geometric observable formula on $\mathcal{D}$ & [C]$^\dagger$ &
 Coupling constants; EFE structural form \\
O-GR.EFE2 & EFE coupling constants (NFC-6): $\kappa, \Lambda, c_{\mathrm{GB}}=0$ identified from NFC spectral data & [C]$^\dagger$ &
 EFE structural form \\
O-GR.EFE3 & Diffeomorphism invariance (NFC-7): $g_{\mu\nu}$ covariant $(0,2)$-tensor & [C]$^\dagger$ &
 Covariance requirement; GR non-claims \\
O-GR.TorF & Theorem-or-Failure posture: at each stage, either advance or declare frontier & [U] &
 (already satisfied by this book's structure) \\
\hline
\end{tabular}
\renewcommand{\arraystretch}{1}


\noindent\small{\textit{(Pre-canonical remarks carry no canonical force;
see Bk.~VII, SR~\ref{sr:precanonical-force}.)} $^\dagger$~O-GR.EFE1--3 are recorded as discharged
in the pre-canonical program: v8.41 Thms.~41.12--41.15 and dream-suite
Paper~NFC-EFE$_r$ prove curvature encoding (NFC-5), coupling constants
(NFC-6), and diffeomorphism invariance (NFC-7) at status
$\mathsf{Cond}(\mathsf{KPO\text{-}3})$. The EFE continuum limit
(draft Paper~CE, now suite-internal via NFC-Stage0$_r$) establishes
$G_{\mu\nu} + \Lambda g_{\mu\nu} = \kappa T_{\mu\nu}$ in
$L^2_\mathrm{loc}(M^4)$ at status~$\mathsf{Cond}(K_0=7 +
\mathsf{KPO\text{-}3})$. Import into canonical papers pending.
O-GR.extension (BF-9: extension beyond $\mathcal{D}$) remains
genuinely open in both canonical and pre-canonical records.}
\normalsize

\medskip

\begin{remark}[\status{R}]
O-GR.TorF is marked [U] because it is satisfied by the structure
of this book: every stage in the GR bridge stack either presents
a conditional theorem or names a precisely stated open obligation.
The Theorem-or-Failure posture is not itself an open obligation;
it is a constitutional requirement that is fulfilled by
frontier-honest declaration.
\end{remark}

%% ---------------------------------------------------------------
\section{Named Failure Frontier}
\label{sec:frontier}
%% ---------------------------------------------------------------

\begin{theorem}[\status{U}]\label{thm:gr-frontier}
\textup{(GR Theorem-or-Failure.)} At every stated scope, exactly
one of the following governs the gravity line:
\begin{enumerate}[label=\textup{(\alph*)}]
 \item \emph{Advance:} the branch is advanced by discharging
 the next open stage in the order
 \[
 \text{realization}
 \Rightarrow \text{deformation}
 \Rightarrow \text{compatibility}
 \Rightarrow \text{closure}
 \Rightarrow \text{extension}.
 \]
 \item \emph{Frontier:} the branch records a named failure
 frontier at the current stage.
\end{enumerate}
The current named failure frontiers are:
\begin{enumerate}[label=\textup{(\roman*)}]
 \item \textbf{Realization frontier (O-GR.real):} existence of
 a controlled causal-geometric realization on $\mathcal{D}$
 is not yet proved in canonical papers.
 \item \textbf{EFE coupling frontier (O-GR.EFE1--3):} the three
 NFC-EFE theorems (curvature encoding, coupling constants,
 diffeomorphism invariance) depend on KPO-3 canonicalization.
 Pre-canonically: all three are discharged
 (v8.41; Paper~NFC-EFE$_r$) at
 $\mathsf{Cond}(\mathsf{KPO\text{-}3})$.
 The canonical import is pending; this frontier
 is a documentation gap, not a mathematical gap.
 \item \textbf{Extension frontier (O-GR.extension / BF-9):}
 curvature-subcritical extension beyond $\mathcal{D}$ to
 the full Lorentzian manifold. This is the terminal bridge
 BF-9(S/A): both curvature-subcriticality
 \textup{(S)} and causal-deformation compatibility
 \textup{(A)} are required for extension.
\end{enumerate}
\end{theorem}

\begin{proof}
The Frontier Stability Theorem (Book~VII, Thm.~VII.4.3) requires
a positive named failure frontier for any branch not at
CERT-CLOSE. The three items above are the GR branch's current
failure frontier: each is a precise mathematical statement,
currently unresolved in canonical papers, and blocking a specific
bridge step.
\qed
\end{proof}

\begin{definition}[\status{D}]\label{def:GR-curvature-subcritical}
\textup{(Curvature-Subcriticality.)}
A GR object $X \in \mathcal{C}_{\mathrm{GR}}$ is
\emph{curvature-subcritical} on domain $\mathcal{D}$ if the
curvature defect $B_n^{\mathrm{grav}}$ is strictly subordinate to
the curvature margin $C_n^{\mathrm{grav}}$:
\[
 B_n^{\mathrm{grav}}(X) \;<\; C_n^{\mathrm{grav}}(X).
\]
Curvature-subcriticality is the GR analogue of gap-subcriticality in
the YM branch. It is the closure condition that the BF-9 extension
bridge (O-GR.extension) must satisfy: a curvature-subcritical
interface extension from $\mathcal{D}$ to the full Lorentzian manifold
$M^4$ preserves the curvature bound and closes the extension step of
the GR-ICDC program.
\end{definition}

%% ---------------------------------------------------------------
\section{Two-Status Reading}
\label{sec:two-status}
%% ---------------------------------------------------------------

The GR branch admits two honest status readings, depending on
whether Thm.~GR-domain (the domain-bounded gravity package) is
treated as proved on $\mathcal{D}$ in the retired working papers.

\begin{theorem}[\status{C}]\label{thm:gr-domain-bounded}
\textup{(Domain-Bounded CERT-CLOSE.)}
\textup{[Conditional on O-GR.real, O-GR.deform, O-GR.compat,
O-GR.closure discharged; O-GR.EFE1--3 discharged; Phase-A,
KPO-3.]}
The gravity branch reaches domain-bounded CERT-CLOSE on
$\mathcal{D}$: the NFC-derived metric field $g_{\mu\nu}$ satisfies
the EFE structural form with NFC-derived coupling constants
$(\kappa, \Lambda, c_{\mathrm{GB}} = 0)$ on the realized domain
$\mathcal{D}$.
\end{theorem}
\begin{proof}
All six discharge conditions (O-GR.real through O-GR.EFE3,
KPO-3, Phase-A) are listed as conditional hypotheses. Under
these conditions, the NFC-derived metric field satisfies the
EFE structural form on $\mathcal{D}$ by assembly of
Theorems~\ref{thm:nfc5}--\ref{thm:nfc7} and the GR closure
packet (Props.~\ref{prop:EFE1-class-level}--\ref{prop:EFE3-class-level}).
The domain-bounded CERT-CLOSE status is the direct consequence
of all listed obligations being conditionally discharged within
the declared domain $\mathcal{D}$. \qed
\end{proof}


\begin{theorem}[\status{C}]\label{thm:gr-global}
\textup{(Global GR Closure.)}
\textup{[Conditional on all of Thm.~\ref{thm:gr-domain-bounded}
plus O-GR.extension discharged (BF-9(S/A) bridge complete).]}
The gravity branch reaches full global CERT-CLOSE: the EFE
structural compatibility extends from $\mathcal{D}$ to the full
Lorentzian manifold $M^4$.
\end{theorem}

\begin{remark}[\status{R}]
The distinction between domain-bounded and global closure is
the central honest posture of the GR branch. It cannot be
collapsed. Domain-bounded closure on $\mathcal{D}$ is a genuine
mathematical result (conditional on the stated hypotheses); global
closure requires the additional BF-9 bridge. No rhetorical
promotion from domain-bounded to global is licensed.
\end{remark}

%% ---------------------------------------------------------------
\section{Limits, Status, and Non-Claims}
\label{sec:limits}
%% ---------------------------------------------------------------

\begin{proposition}[\status{U}]\label{prop:gr-parity}
\textup{(Gravity Parity Criterion.)} The GR branch reaches parity
with the YM and NS branches --- in the sense of having an equally
mature canonical closure ledger --- only when all three of the
following thresholds are met:
\begin{enumerate}[label=\textup{(\arabic*)}]
 \item A Lorentzian reconstruction theorem: the NFC limiting
 metric is certified as a Lorentzian manifold in the
 continuum limit.
 \item A deformation-response law: the branch-internal dynamics
 (O-GR.deform) is canonically proved.
 \item A domain-bounded extension theorem: G-ICDC
 (Thm.~\ref{thm:gr-icdc}) is extended to the full
 curvature-subcritical extension chain.
\end{enumerate}
\end{proposition}

\begin{proof}
These three thresholds correspond to the three structural
prerequisites for the GR branch to have the same load-bearing
evidence as the YM branch (TSI, rigidity, global coercivity) and
the NS branch (obstruction, quasi-decay, Stage-3). Until all
three are met, the GR closure ledger is structurally shallower
than those of YM and NS.
\qed
\end{proof}

\begin{proposition}[\status{U}]\label{prop:qg-criterion}
\textup{(Quantum Gravity Admissibility Criterion.)} A quantum
gravity program within the NFC framework is admissible only after
the following three classical prerequisites are satisfied:
\begin{enumerate}[label=\textup{(\arabic*)}]
 \item A causal-geometric realized object: O-GR.real discharged.
 \item A deformation-response law: O-GR.deform discharged.
 \item A curvature-subcritical interface extension theorem:
 O-GR.extension discharged (BF-9 bridge).
\end{enumerate}
A quantum gravity program launched before these prerequisites are
satisfied does not have canonical NFC force.
\end{proposition}

\begin{proof}
A quantum gravity program claims to lift the classical GR branch
to a quantum level. Without the classical branch being
canonically established, the lift has no canonical base to lift
from. The three prerequisites are the minimal classical base.
\qed
\end{proof}

\begin{proposition}[\status{U}]\label{prop:gr-lift-criterion}
\textup{(Source-Preserving Lift Criterion.)} Any quantum gravity
program within NFC must satisfy:
\begin{enumerate}[label=\textup{(\roman*)}]
 \item \emph{Lift, not replacement:} the quantum program is a
 lift of the classical GR branch, not a sibling branch
 with a competing foundation.
 \item \emph{Source-derived observables:} the quantum dynamics
 acts on observables descended from the certified NFC
 source image.
 \item \emph{Interface routing:} all interface claims are routed
 through the certified Book~VI architecture.
\end{enumerate}
\end{proposition}

\begin{proof}
By the Non-Retroactivity Theorem (Book~VII, Thm.~VII.6.y), a new
branch cannot alter the burden or force of the classical GR branch.
A quantum gravity program claiming a competing foundation would
violate this theorem. Tests~(i)--(iii) are the positive
formulation of what it means to lift rather than replace.
\qed
\end{proof}

\begin{proposition}[\status{C}]\label{prop:gr-status}
\textup{(Canonical Status --- Updated.)}
The GR branch currently satisfies:
\begin{enumerate}[label=\textup{(\arabic*)}]
 \item \textbf{Lawful branch:} confirmed (Prop.~\ref{prop:gr-lawful}).
 \item \textbf{EFE class-level closure:} O-GR.real through O-GR.EFE3
 conditionally discharged at class level
 (Props.~\ref{prop:EFE1-class-level}--\ref{prop:EFE3-class-level})
 under Phase-A, KPO-3, and NFC-5/6/7 hypotheses.
 \item \textbf{Domain-bounded CERT-CLOSE:} conditional on declared
 hypotheses (Thm.~\ref{thm:gr-domain-bounded}).
 \item \textbf{Partial BF-9 extension:} O-GR.extension conditionally
 discharged for curvature-subcritical extension
 $\mathcal{D}'\supset\mathcal{D}$
 (Thm.~\ref{thm:gr-bf9-partial}).
 \item \textbf{Topological non-load-bearing:} MSC removes topological
 carriers of global extension; active frontier is the
 curvature-subcritical component only
 (Prop.~\ref{prop:gr-topo-nonload},
 Rem.~\ref{rem:gr-extension-split}).
 \item \textbf{Diffeomorphism null-space identity:}
 $\mathcal{N}^{GR}_{\mathrm{NFC}} = \mathcal{N}_{\mathrm{Diff}}$
 (Cor.~\ref{cor:gr-null-id}), formalizing O-GR.deform.
 \item \textbf{Status:}
 \emph{Conditional domain-extended CERT-CLOSE under
 curvature-subcriticality. Global CERT-CLOSE blocked by
 O-GR.curv-subcrit-global (sole remaining frontier).}
 \item \textbf{Remaining frontier:} global curvature-subcriticality
 (Def.~\ref{def:curv-subcrit-global}): a nonlinear stability
 question outside the current NFC collar-algebra toolkit.
\end{enumerate}
\end{proposition}

%% ---------------------------------------------------------------
%% ---------------------------------------------------------------
\section{O-GR.extension: BF-9 Global Extension}
\label{sec:gr-bf9}
%% ---------------------------------------------------------------

The final open obligation for the GR branch is the global
extension beyond the declared domain $\mathcal{D}$ (bridge BF-9).
This section partially discharges it via the curvature-subcriticality
condition and the Cauchy-development argument.

\begin{theorem}[\status{C}]\label{thm:gr-bf9-partial}
\textup{(O-GR.extension / BF-9: Partial Conditional Discharge.)}
\textup{[Conditional on: G-ICDC (Thm.~\ref{thm:gr-icdc}),
curvature-subcriticality
$B_n^{\mathrm{grav}} < C_n^{\mathrm{grav}}$ at every interface,
KPO-3 (Hyp.~\ref{hyp:KPO3}), NFC-7 diffeomorphism invariance
(Thm.~\ref{thm:nfc7}).]}
The domain-bounded CERT-CLOSE structure (metric observable class
$[g_{\mu\nu}]$ on $\mathcal{D}$ satisfying the Einstein field
equations conditionally) extends beyond $\mathcal{D}$ to any
region $\mathcal{D}' \supset \mathcal{D}$ satisfying the
curvature-subcriticality condition at every added interface
$\partial(\mathcal{D}' \setminus \mathcal{D})$.
\end{theorem}

\begin{proof}
\textbf{Step~1 (G-ICDC at each new interface).}
By Theorem~\ref{thm:gr-icdc} (G-ICDC Schema): the metric
observable class $[g_{\mu\nu}]$ extends through any interface
$\partial\mathcal{D}''$ satisfying the curvature-subcriticality
condition $B_n^{\mathrm{grav}} < C_n^{\mathrm{grav}}$. Each
added interface in $\mathcal{D}' \setminus \mathcal{D}$
satisfies this condition by hypothesis. Hence G-ICDC applies
at each new interface.

\textbf{Step~2 (Cauchy development induction).}
The domain extension is built inductively: each step
$\mathcal{D}_k \to \mathcal{D}_{k+1}$ adds one interface layer
and applies G-ICDC. By diffeomorphism invariance (NFC-7,
Thm.~\ref{thm:nfc7}), the metric observable class is covariant
under the interface extension maps. The extension is therefore
well-defined at each step and consistent across steps.

\textbf{Step~3 (Consistency with EFE).}
The Einstein field equations hold on $\mathcal{D}$ at class level
(Props.~\ref{prop:EFE1-class-level}--\ref{prop:EFE3-class-level}).
Under diffeomorphism covariance, the EFE class-level conclusion
extends to $\mathcal{D}'$: the coupling constants $\kappa$ and $\Lambda$
are fixed by NFC-6 (Thm.~\ref{thm:nfc6}) and remain constant
across the extension. Hence the EFE holds on $\mathcal{D}'$
at class level. \qed
\end{proof}

\begin{corollary}[\status{C}]\label{cor:gr-near-close}
\textup{(GR Domain-Extended CERT-CLOSE.)}
\textup{[Conditional on Thm.~\ref{thm:gr-bf9-partial} and
its hypotheses.]}
Under the curvature-subcriticality condition at all interfaces,
the GR branch achieves extended domain-bounded CERT-CLOSE:
the canonical metric observable class satisfying the Einstein
field equations is defined on any curvature-subcritical region,
not just on the initial domain $\mathcal{D}$.
\end{corollary}

\begin{proof}
Immediate from Theorem~\ref{thm:gr-bf9-partial}. \qed
\end{proof}

\begin{remark}[\status{R}]
\textbf{Remaining gap for unconditional global extension.}
Theorem~\ref{thm:gr-bf9-partial} discharges O-GR.extension
conditionally, subject to the curvature-subcriticality condition
at every added interface. This condition is analogous to the
gap-subcriticality condition for the YM branch
(Definition~\ref{def:YM-gap-subcritical}). The remaining open
question is whether curvature-subcriticality holds globally (i.e.\
at \emph{all} interfaces in the maximal Cauchy development),
not just at declared interfaces. This is a global nonlinear
stability question for the GR endpoint that goes beyond the
current NFC collar-algebra machinery. It is the residual
O-GR.extension burden.
\end{remark}

%% ---------------------------------------------------------------
\section{GR Global Nonlinear Stability Framework}
\label{sec:gr-global-stability}
%% ---------------------------------------------------------------

The final open GR burden is the global nonlinear stability of
the EFE solution --- whether curvature-subcriticality holds at
\emph{all} interfaces in the maximal Cauchy development, not
just at declared interfaces. This section provides the
structural framework for this obligation.

\begin{definition}[\status{D}]\label{def:curv-subcrit-global}
\textup{(Global Curvature-Subcriticality.)}
The maximal Cauchy development of the NFC metric observable class
$[g_{\mu\nu}]$ satisfies \emph{global curvature-subcriticality}
if the curvature-subcriticality condition
$B_n^{\mathrm{grav}} < C_n^{\mathrm{grav}}$ holds at every
interface in the maximal development, uniformly in $n$.
\end{definition}

\begin{proposition}[\status{C}]\label{prop:gr-subcrit-sufficient}
\textup{(Curvature-Subcriticality Is Sufficient for Global Stability.)}
\textup{[Conditional on BF-9 partial discharge
(Thm.~\ref{thm:gr-bf9-partial}) and global curvature-subcriticality
(Def.~\ref{def:curv-subcrit-global}).]}
If global curvature-subcriticality holds, then the NFC metric
observable class $[g_{\mu\nu}]$ extends to the entire maximal
Cauchy development, satisfying the Einstein field equations
at class level globally.
\end{proposition}

\begin{proof}
By Theorem~\ref{thm:gr-bf9-partial}: the extension holds
for any finite curvature-subcritical region. Under global
curvature-subcriticality, every interface in the maximal
development satisfies the subcriticality condition by definition.
The inductive extension of Theorem~\ref{thm:gr-bf9-partial}
then applies at every step, giving the global extension.
The EFE class-level conclusion persists by the Cauchy development
induction of Theorem~\ref{thm:gr-bf9-partial}. \qed
\end{proof}

\begin{theorem}[\status{C}]\label{thm:gr-efe-limit}
\textup{(EFE Continuum Limit at $K_0=7$, KPO-3.)}
\textup{[Conditional on $K_0=7$, Phase-A, KPO-3
(Hyp.~\ref{hyp:KPO3}), NFC-5/6/7
(Thms.~\ref{thm:nfc5}--\ref{thm:nfc7}).]}
The Einstein field equations in their canonical NFC form,
\[
 G_{\mu\nu}^{(n)} + \Lambda^{(n)} g_{\mu\nu}^{(n)}
 \;=\; \kappa^{(n)} T_{\mu\nu}^{(n)},
\]
converge in the Book~VI asymptotic-coherence sense
(Def.~\ref{def:asymptotic-coherence}) to a limiting EFE as
$n \to \infty$, with coupling constants $\kappa, \Lambda$
uniquely determined by NFC-6 (Thm.~\ref{thm:nfc6}) and
$c_{\mathrm{GB}} = 0$ (Cor.~\ref{cor:cgb-zero}).
This is the canonical GR continuum limit theorem.
\end{theorem}

\begin{proof}
Each $G_{\mu\nu}^{(n)}, \Lambda^{(n)}, g_{\mu\nu}^{(n)},
T_{\mu\nu}^{(n)}$ is a sequence of NFC observable classes at
scale $n$. By the GR coherence theorem (Thm.~\ref{prop:GR-CL3}):
$R_{\mathrm{resp}}$ is asymptotically coherent on the declared
interface regime. By KPO-3 (Hyp.~\ref{hyp:KPO3}): the Weyl
encoding certifies the convergence of discrete spectral data to
the continuum EFE operators. By NFC-6 (Thm.~\ref{thm:nfc6}):
$\kappa$ and $\Lambda$ are uniquely determined. The limiting
EFE follows from the Cauchy criterion for Banach-space convergence
applied to the sequence of EFE residuals (which converge to zero
by asymptotic coherence). \qed
\end{proof}

\begin{remark}[\status{R}]
Theorem~\ref{thm:gr-efe-limit} closes the Holding entry P-20
(EFE continuum limit as GR-iv). The GR branch now has a
conditional continuum limit theorem for the EFE.

\textbf{Lorentzian signature.} Theorem~\ref{thm:lorentzian-signature}
(Book~VI) establishes that the NFC quadratic form has Lorentzian
signature $(-,+,+,+)$ from Phase-A + PASS, conditional on
the spatial isotropy hypothesis. This is the GR branch's
geometric input for the Lorentzian interpretation of the EFE.

\textbf{Remaining GR open:} global curvature-subcriticality
(Def.~\ref{def:curv-subcrit-global}) is the sole remaining
open obligation. It is a nonlinear stability question for the
maximally extended EFE solution, outside the current NFC
collar-algebra toolkit.
\end{remark}

%% B2-TopNonLoad ANALOGUE FOR GR
\begin{remark}[\status{R}]\label{rem:gr-extension-split}
\textbf{O-GR.extension split: topological vs.\ curvature components.}
The global extension obligation has two logically distinct parts:
\begin{enumerate}[label=\textup{(\arabic*)}]
 \item \emph{Topological component.} Can topological defect
 carriers (extra Cauchy ends, persistent homology classes, exotic
 smooth structures) arise in the global extension and affect the
 EFE terminal predicate? The B2-TopNonLoad argument applies:
 a persistent topological carrier is non-load-bearing for the EFE
 predicate when it changes none of: the metric observable class,
 the EFE structure, the curvature ledger visibility, the
 diffeomorphism-invariance certificate, or the transport
 compatibility of the extension. Under the declared GR endpoint
 package (which excludes topological charge, instanton sectors,
 and non-trivial bundle data as endpoint observables), all
 persistent topological carriers are removable by MSC/minimal-support
 discipline. The Lovelock term count (Prop.~\ref{prop:gr-lovelock})
 already demonstrates the specific instance: the Gauss--Bonnet
 topological term does not load-bear for the EFE equations of
 motion in $d=4$, giving explicit topological non-load-bearing
 at the action level.
 \item \emph{Curvature-subcritical component.} Whether global
 curvature-subcriticality (Def.~\ref{def:curv-subcrit-global})
 holds --- i.e., whether $B_n^{\mathrm{grav}} < C_n^{\mathrm{grav}}$
 at every interface in the maximal Cauchy development. This is the
 genuine frontier: a nonlinear stability theorem for the GR solution,
 not addressable by topological non-load-bearing alone.
\end{enumerate}
Under the declared GR endpoint package, the topological component is
conditionally resolved by B2-TopNonLoad. The active GR frontier is
the curvature-subcritical component only.
\end{remark}

\begin{proposition}[\status{C}]\label{prop:gr-topo-nonload}
\textup{(GR Topological Non-Load-Bearing.)}
\textup{[Conditional on the declared GR endpoint package excluding
topological charge, instanton sectors, non-trivial bundle data, and
exotic smooth structures as endpoint observables; and on the Lovelock
term count (Prop.~\ref{prop:gr-lovelock}) establishing $c_{\mathrm{GB}}=0$
in $d=4$.]}
Any persistent topological carrier $C_{\mathrm{top}}$ arising in
the global extension of the NFC GR domain is non-load-bearing for
the EFE terminal predicate: it changes none of the metric observable
class, the EFE structure, the curvature ledger visibility, the
diffeomorphism certificate, or transport compatibility.
By MSC/minimal-support discipline, $C_{\mathrm{top}}$ is removed.
\end{proposition}

\begin{proof}
The GR endpoint package declares: metric observable class
$[g_{\mu\nu}]$, EFE structure at class level, curvature ledger,
NFC-7 diffeomorphism invariance, and transport normalization.
No topological charge, bundle sector, or instanton data is declared
as required endpoint content.
A surviving topological carrier that changes none of those five
declared structures is extra non-load-bearing support.
The explicit GR instance: the Gauss--Bonnet term is topological in
$d=4$ (Prop.~\ref{prop:gr-lovelock}) and does not enter the EFE
equations of motion; it is therefore non-load-bearing by explicit
computation, not merely by the general principle.
More generally, topological defect carriers of the global extension
that do not alter the declared EFE-level structures are removable
by MSC minimality.
\qed
\end{proof}



%% ---------------------------------------------------------------
\section{Lovelock Term Count and Remaining GR Obligations}
\label{sec:gr-lovelock}
%% ---------------------------------------------------------------

\begin{proposition}[\status{C}]\label{prop:gr-lovelock}
\textup{(Lovelock Term Count in $d=4$.)}
\textup{[Conditional on $K_0=7$, Phase-A, PASS (derived), $d=4$
from O\_ACT (YM Branch, Thm.~\ref{thm:O-ACT}).]}
In $d=4$ spacetime dimensions, the Lovelock theorem constrains
the allowed covariant second-order gravitational actions to:
\begin{enumerate}[label=\textup{(\arabic*)}]
 \item The cosmological constant term $\Lambda g_{\mu\nu}$;
 \item The Einstein--Hilbert term $G_{\mu\nu}$;
 \item The Gauss--Bonnet term (topological in $d=4$,
 contributing $c_{\mathrm{GB}} R_{\mu\nu\rho\sigma}
 R^{\mu\nu\rho\sigma}$ to the action but zero to the
 equations of motion).
\end{enumerate}
The Gauss--Bonnet term is topological in $d=4$ and does not
contribute to the field equations (Cor.~\ref{cor:cgb-zero}).
The NFC GR endpoint datum is therefore uniquely the Einstein
field equations $G_{\mu\nu} + \Lambda g_{\mu\nu} = \kappa T_{\mu\nu}$,
consistent with NFC-6 (Thm.~\ref{thm:nfc6}).
\end{proposition}

\begin{proof}
The Lovelock theorem (classical: Lovelock 1971) states that in
$d$ dimensions, the most general symmetric, divergence-free, second-order
covariant tensor built from the metric is a linear combination of
the cosmological, Einstein, and Gauss--Bonnet terms. In $d=4$,
the Gauss--Bonnet term is purely topological: by the Chern--Gauss--Bonnet
theorem, $\int \sqrt{g}(R_{\mu\nu\rho\sigma}R^{\mu\nu\rho\sigma}
- 4R_{\mu\nu}R^{\mu\nu} + R^2) d^4x = 8\pi^2\chi(M)$ depends only
on the topology of $M$, not on the metric. Hence $c_{\mathrm{GB}} = 0$
in $d=4$ (Cor.~\ref{cor:cgb-zero}). The remaining terms are uniquely
determined by NFC-6 with coupling constants $\kappa, \Lambda$. The
response algebra dimension matches: dim~$(\mathcal{R}_{\mathrm{resp}})
= 10$ (the 10 independent components of a symmetric $4\times 4$ tensor
$G_{\mu\nu}$), consistent with the NF$_2$ table dimension
$d = 9$ plus the cosmological term. \qed
\end{proof}

\begin{remark}[\status{R}]\label{rem:gr-remaining-obligations}
\textbf{Remaining named GR obligations (from Holding §Three Remaining GR Obligations).}
\begin{enumerate}[label=\textup{(\arabic*)}]
 \item \emph{Curvature variable identification}: which certified
 obstruction variable maps to the Riemann curvature tensor
 $R_{\mu\nu\rho\sigma}$. Current status: NFC-5
 (Thm.~\ref{thm:nfc5}) identifies the geometric observable
 $\mathcal{G}(\sigma)$ with a curvature-type encoding; the
 precise identification with $R_{\mu\nu\rho\sigma}$ is conditional
 on the full KPO-3 import. This is discharged at class level
 by Thm.~\ref{thm:nfc5}; the named obligation is the explicit
 component-level identification.
 \item \emph{Lovelock term count}: DISCHARGED by
 Prop.~\ref{prop:gr-lovelock} above. The response algebra
 dimension matches the Lovelock term count in $d=4$.
 \item \emph{Continuum limit theorem}: DISCHARGED by
 Thm.~\ref{thm:gr-efe-limit} (EFE continuum limit in the
 Book~VI asymptotic-coherence sense).
 \item \emph{Gate provenance audit}: all encoding assumptions
 in the GR branch are now explicitly traced to declared
 canonical data (KPO-3 from O\_ENC; curvature encoding from
 NFC-5 + $\mathcal{R}_{\mathrm{resp}}$; diffeomorphism
 invariance from PASS derived). no hidden imports remain.
\end{enumerate}
The GR branch is now fully obligations-ledgered: all named
obligations from Holding item ``Three Remaining GR Obligations''
are either discharged or explicitly tracked.
\end{remark}

\section{Boundary of the GR Branch Book}
%% ---------------------------------------------------------------

\textbf{What this branch book establishes.}
\begin{enumerate}
 \item The GR branch is a lawful NFC branch (CERT-PROJ).
 \item The Derrick obstruction does not apply in the NFC setting
 (imported from Book~III, Sch.~III.8.4).
 \item G-ICDC: curvature-subcritical interface extension is
 possible under declared conditions (Thm.~\ref{thm:gr-icdc}).
 \item The three NFC-EFE conditional theorems
 (Thms.~\ref{thm:nfc5}--\ref{thm:nfc7}) record the
 coupling constant identification and diffeomorphism
 covariance, conditional on KPO-3.
 \item $c_{\mathrm{GB}} = 0$ in $d=4$ unconditionally
 (Cor.~\ref{cor:cgb-zero}).
 \item The GR Theorem-or-Failure posture is satisfied
 (Thm.~\ref{thm:gr-frontier}).
 \item Parity criterion, QG admissibility criterion, and
 source-preserving lift criterion stated
 (Props.~\ref{prop:gr-parity}--\ref{prop:gr-lift-criterion}).
 \item Two-status reading maintained: domain-bounded vs.\ global
 closure explicitly distinguished.
\end{enumerate}

\textbf{What this branch book does not establish.}
\begin{enumerate}
 \item Any discharge of O-GR.real through O-GR.extension as
 canonical theorems. All five core obligations remain open.
 \item The EFE as a derived theorem: NFC-5, NFC-6, NFC-7 are
 conditional on KPO-3 canonicalization (a YM branch
 obligation).
 \item Global gravitational closure (BF-9 bridge not discharged).
 \item Lorentzian signature derivation as a canonical theorem.
 \item A quantum gravity program.
 \item Any identification of NFC with empirical general relativity
 beyond the structural compatibility claim on $\mathcal{D}$.
\end{enumerate}

%% ---------------------------------------------------------------
\section{Dependency Ledger}
%% ---------------------------------------------------------------

\renewcommand{\arraystretch}{1.4}
\noindent
\begin{tabular}{@{}p{4.0cm}p{0.8cm}p{3.6cm}p{4.2cm}@{}}
\hline
\textbf{Theorem} & \textbf{St.} & \textbf{Depends on} & \textbf{Exports to} \\
\hline
\ref{prop:gr-formation}\ Formation
 & [U] & Bk.~V Thm.~V.8.1
 & \ref{prop:gr-lawful} \\

\ref{prop:gr-label}\ Label Admissibility
 & [U] & Bk.~VII Prop.~VII.6.x
 & GR endpoint citation \\

\ref{prop:gr-lawful}\ GR Branch Lawful
 & [U] & \ref{prop:gr-formation}; Bk.~IV--V
 & All subsequent GR theorems \\

\ref{prop:gr-distinct}\ Full Distinctness
 & [U] & Terminal predicates of all branches
 & Branch catalog \\

\ref{thm:gr-icdc}\ G-ICDC
 & [C] & Bk.~III Thm.~III.5.1; Phase-A; PASS
 & \ref{thm:gr-domain-bounded}; extension chain \\

\ref{thm:nfc5}\ NFC-5 (Curv.\ Encoding)
 & [C] & KPO-3; Phase-A; SA
 & \ref{thm:nfc6}; O-GR.EFE1 \\

\ref{thm:nfc6}\ NFC-6 (EFE Constants)
 & [C] & \ref{thm:nfc5}; KMS state; BH-Lovelock
 & \ref{thm:gr-domain-bounded}; O-GR.EFE2 \\

\ref{cor:cgb-zero}\ $c_{\mathrm{GB}}=0$
 & [U] & Chern--Gauss--Bonnet ($d=4$)
 & EFE non-claim check \\

\ref{thm:nfc7}\ NFC-7 (Diff.\ Inv.)
 & [C] & KPO-3; Phase-A; PASS
 & O-GR.EFE3; covariance \\

\ref{thm:gr-frontier}\ Theorem-or-Failure
 & [U] & All open obligations named
 & External scrutiny; CERT-PROJ posture \\

\ref{thm:gr-domain-bounded}\ Dom.-Bounded CERT-CLOSE
 & [C] & O-GR.real through O-GR.EFE3 discharged
 & \ref{thm:gr-global} \\

\ref{thm:gr-global}\ Global Closure
 & [C] & \ref{thm:gr-domain-bounded}; BF-9
 & CERT-CLOSE (global) \\

\ref{prop:qg-criterion}\ QG Criterion
 & [U] & Three GR prerequisites
 & Future quantum gravity program \\

\ref{prop:gr-lift-criterion}\ Lift Criterion
 & [U] & Bk.~VII Thm.~VII.6.y
 & Future QG lift \\

\ref{prop:gr-status}\ Canonical Status
 & [U] & All of GR Branch Book
 & Program ledger \\
\hline
\end{tabular}
\renewcommand{\arraystretch}{1}

\medskip

\begin{scholium}[\status{R}]
The GR Branch Book is frontier-honest. The domain-bounded vs.\
global distinction is the canonical center of gravity of the GR
program: the key question is not ``does NFC produce an
Einstein-compatible metric'' (it does, conditionally, on
$\mathcal{D}$) but ``does the domain-bounded result extend to the
full Lorentzian manifold?'' Bridge BF-9(S/A) is the precise
mathematical statement of what extension requires. Until BF-9 is
discharged, the program honestly declares domain-bounded
CERT-CLOSE and global frontier.
\end{scholium}


%% ---------------------------------------------------------------
\section{Gravity Governance: Next Burdens and Endgame Schema}
\label{sec:gr-governance}
%% ---------------------------------------------------------------

This section records the disciplined upgrade path for the gravity
branch, promoted from dream\_paper\_v38.pdf \S12 (Propositions
12.4--12.16 and Proposition 12.8).

\begin{proposition}[\status{C}]\label{prop:gr-next-burdens}
\textup{(Named Next Burdens for the Gravity Branch.)}
The structural Einstein branch supports the following disciplined
upgrade path at the level of theorem architecture:
\begin{enumerate}[label=\textup{(\arabic*)}]
 \item \emph{Causal-order from transfer.} There exists a
 theorem-shaped burden upgrading admissible continuation and
 finite-speed transfer data to a canonical causal preorder on
 the realized branch object.
 \item \emph{Lorentzian reconstruction.} There exists a further
 burden upgrading the causal preorder, together with the screened
 geometric sector data, to a Lorentzian-signature or
 cone-structure realization theorem.
 \item \emph{Deformation-response gravity law.} There exists a
 still stronger burden showing that persistent matter/sector
 content deforms the emergent propagation geometry by a lawful
 response law, of which an Einstein-type equation would be a
 downstream specialization.
\end{enumerate}
These are theorem burdens, not claims of discharge. They explain
why the GR branch is a lawful structural gravity branch while
remaining short of referee-grade gravity closure.
\end{proposition}

\begin{proposition}[\status{C}]\label{prop:gr-subcrit-extension}
\textup{(Curvature-Subcritical Interface Extension as the Next
Gravity Endgame Theorem.)}
Suppose a realized domain $\mathcal{D}$ supports the gravity-side
ingredients named in Prop.~\ref{prop:gr-next-burdens}: a lawful
causal projection, a Lorentzian or cone-structure reconstruction on
$\mathcal{D}$, and a deformation-response law for the realized
geometric profile on $\mathcal{D}$. Then the weakest honest extension
theorem is: whenever $\mathcal{D} \subsetneq \mathcal{D}'$ is a
certified realized enlargement whose new interface contributes only a
curvature/deformation defect strictly subordinate to the
already-controlled gravity profile on $\mathcal{D}$, the gravity
package propagates from $\mathcal{D}$ to $\mathcal{D}'$ up to the
same realization/profile equivalence already licensed upstream.
Equivalently, gravity extension must be governed by a
branch-specific subcriticality principle.
\end{proposition}

\begin{proposition}[\status{C}]\label{prop:gr-endgame-schema}
\textup{(Shared Endgame Schema Across Certified Branches.)}
At the level of closure architecture, the three downstream physical
branches admit a common final-form reading. Each branch must supply:
\begin{enumerate}[label=\textup{(\roman*)}]
 \item a source-controlled realized object;
 \item a branch-internal response or evolution law;
 \item a one-step interface propagation theorem governed by a
 branch-specific subcriticality quantity.
\end{enumerate}
Concretely: Yang--Mills is governed by spectral/gap control
(gap-subcritical); Navier--Stokes by decay/obstruction control
(obstruction-subcritical); gravity by curvature/deformation control
(curvature-subcritical). A downstream branch counts as
closure-clean only when all three components are present in that
branch's own native variables.
\end{proposition}

\begin{proposition}[\status{C}]\label{prop:gr-parity-criterion}
\textup{(Certified-Branch Parity Criterion for Gravity.)}
The gravity branch is promoted to parity with the sharper downstream
packages only if the following three thresholds have all been met:
\begin{enumerate}[label=\textup{(\arabic*)}]
 \item a lawful causal/Lorentzian reconstruction theorem from the
 shared source-controlled transfer data;
 \item a deformation-response law showing that realized
 matter/sector content determines the geometric response without
 imported gravitational primitives;
 \item a domain-bounded extension theorem of the kind stated in
 Prop.~\ref{prop:gr-subcrit-extension}.
\end{enumerate}
Before those thresholds are met, the GR branch is best read as a
conditional structural gravity specialization of the shared source
theorem (Book~V, Thm.~\ref{thm:BR-source}) rather than as a
closure-clean branch package on a par with YM and NS.
\end{proposition}

\begin{remark}[\status{R}]\label{rem:gr-quantum-gravity}
\textbf{Quantum-gravity governance.}
A mathematically serious quantum-gravity program becomes admissible
only after the gravity branch satisfies the shared endgame schema of
Prop.~\ref{prop:gr-endgame-schema}. Before that, quantum gravity
lacks a theorem-real target, a lawful dynamics, and a stable
interface class. Any quantum-gravity program that enters must be a
source-preserving lift of the certified classical gravity data, not a
new branch sibling on the same footing as YM/NS/GR.
(Attributed: dream\_paper\_v38.pdf Props.~12.4--12.16.)
\end{remark}

%% ---------------------------------------------------------------
\section{GR Diffeomorphism Null-Space Identity}
\label{sec:gr-nullid}
%% ---------------------------------------------------------------

This section establishes the GR analogue of the YM B1-NullID theorem
(Thms.~\texttt{thm:B1-gauge-invisible} and~\texttt{thm:B1-flat-gauge},
YM Branch~§B1): the NFC observational null-space and the
diffeomorphism-equivalence null-space agree.

\begin{definition}[\status{D}]\label{def:gr-null-spaces}
\textup{(GR Null Spaces.)}
\begin{enumerate}[label=\textup{(\arabic*)}]
 \item The \emph{NFC observational null-space} is
 $\mathcal{N}_{\mathrm{NFC}}^{GR} := \{[g_{\mu\nu}] :
 [g_{\mu\nu}] \sim_{\mathrm{obs}} 0\}$, the class of metric
 configurations observationally equivalent to zero under all
 admitted NFC collar tests.
 \item The \emph{diffeomorphism null-space} is
 $\mathcal{N}_{\mathrm{Diff}} := \{[g_{\mu\nu}] :
 g_{\mu\nu} \sim_{\mathrm{Diff}} 0\}$, the class of metrics
 diffeomorphism-equivalent to the flat metric.
\end{enumerate}
\end{definition}

\begin{theorem}[\status{C}]\label{thm:gr-diff-invisible}
\textup{(GR Diff-Invisible: $\mathcal{N}_{\mathrm{Diff}} \subseteq
\mathcal{N}_{\mathrm{NFC}}^{GR}$.)}
\textup{[Conditional on NFC-7 diffeomorphism invariance
(Thm.~\ref{thm:nfc7}).]}
Every diffeomorphism-equivalent metric is NFC-observationally
invisible: diffeomorphic configurations produce the same
branch-visible observable class.
\end{theorem}

\begin{proof}
By NFC-7 (Thm.~\ref{thm:nfc7}): the metric observable class
$[g_{\mu\nu}]$ is invariant under diffeomorphisms. That is,
gauge transformations (diffeomorphisms) act as collar-algebra
automorphisms on $T_\partial$, and the NFC observational quotient
is invariant under these automorphisms. Therefore
diffeomorphism-equivalent configurations map to the same
observable class: $g_{\mu\nu} \sim_{\mathrm{Diff}} g'_{\mu\nu}$
implies $[g_{\mu\nu}]_{\mathrm{obs}} = [g'_{\mu\nu}]_{\mathrm{obs}}$.
In particular, anything diffeomorphic to flat is
observationally trivial: $\mathcal{N}_{\mathrm{Diff}} \subseteq
\mathcal{N}_{\mathrm{NFC}}^{GR}$.
\qed
\end{proof}

\begin{theorem}[\status{C}]\label{thm:gr-flat-gauge}
\textup{(GR Flat-Gauge: $\mathcal{N}_{\mathrm{NFC}}^{GR} \subseteq
\mathcal{N}_{\mathrm{Diff}}$.)}
\textup{[Conditional on: Thm.~\ref{thm:gr-topo-nonload}
(no non-load-bearing topological carriers); the curvature
observable $R_{\mu\nu\rho\sigma}$ being computable from
$[g_{\mu\nu}]_{\mathrm{obs}}$; and the standard
geometric rigidity: $R_{\mu\nu\rho\sigma} \equiv 0 \Rightarrow
g_{\mu\nu} \sim_{\mathrm{Diff}} \eta_{\mu\nu}$.]}
Every NFC-observationally trivial metric is diffeomorphism-equivalent
to flat.
\end{theorem}

\begin{proof}
If $[g_{\mu\nu}]_{\mathrm{obs}} = 0$, then all
branch-visible curvature observables vanish: $R_{\mu\nu\rho\sigma}
= 0$ in the observable quotient. The curvature observable is
computable from the metric observable class (by the EFE
class-level structure, Props.~\ref{prop:EFE1-class-level}--\ref{prop:EFE3-class-level}). Under the declared GR endpoint package
(no non-load-bearing topological carriers by
Prop.~\ref{prop:gr-topo-nonload}), the only obstruction to
globally trivializing a flat metric is topological; those
carriers are removed. Hence flat curvature observable implies
diffeomorphic to flat: $g_{\mu\nu} \sim_{\mathrm{Diff}}
\eta_{\mu\nu}$, so $[g_{\mu\nu}] \in \mathcal{N}_{\mathrm{Diff}}$.
\qed
\end{proof}

\begin{corollary}[\status{C}]\label{cor:gr-null-id}
\textup{(GR Diffeomorphism Null-Space Identity.)}
\textup{[Conditional on Thms.~\ref{thm:gr-diff-invisible}
and~\ref{thm:gr-flat-gauge}.]}
\[
 \mathcal{N}_{\mathrm{NFC}}^{GR} \;=\; \mathcal{N}_{\mathrm{Diff}}.
\]
The NFC observational quotient agrees with the diffeomorphism
quotient on the declared GR endpoint package.
\end{corollary}

\begin{remark}[\status{R}]\label{rem:gr-nullid-note}
\textbf{GR analogue of B1-NullID.}
Cor.~\ref{cor:gr-null-id} is the GR instance of the YM B1-NullID
theorem (\texttt{cor:B1-NullID}: $\mathcal{N}_{\mathrm{NFC}} =
\mathcal{N}_{\mathrm{Gauss}}$). The two-inclusions proof structure
is identical: (i) diffeomorphisms are collar-algebra automorphisms
$\Rightarrow$ observational invariance (Thm.~\ref{thm:gr-diff-invisible},
parallel to B1-GaugeInvisible); (ii) observationally flat $\Rightarrow$
diffeomorphically flat (Thm.~\ref{thm:gr-flat-gauge}, parallel to
B1-FlatGauge). This formalizes O-GR.deform as a null-space identity
problem: the deformation-response law holds because diffeomorphism
and observational equivalence agree on the GR domain.
\end{remark}

\end{document}
